\section{The Darkness of Excess}

The success of the ascent creates its own danger. The intellect has concluded something true: there is one unreceived source, pure act and subsistent being, simple, unbounded, and intimately causal to all that is. It may now mistake true predication for possession. It may turn ``God'' into a concept small enough to survey.

The patristic and mystical tradition refuses that reduction. St.~Gregory Nazianzen distinguishes conviction that God is from comprehension of what God is. Created effects can lead the mind to their efficient and sustaining cause, yet the divine nature remains beyond enclosure by a created intellect (\emph{Oration} 28.4--6). Elsewhere he gathers biblical and philosophical language into the image of a limitless ``Sea of Being,'' glimpsed only partially, drawing the mind by what it reveals and awakening wonder by what exceeds it (\emph{Oration} 45.3).

This is not an embarrassment added after the proof. It follows from the proof. If the divine essence were exhaustively comprehensible by a finite intellect, the unlimited would be contained by the limited. If God belonged to a genus, stood under principles more universal than himself, or possessed existence according to a limited mode, the regress would not have reached its source. The same reasoning that affirms God therefore requires the mind to deny of him every creaturely limitation.

Aquinas announces this discipline at the entrance to his treatment of the divine essence. Since we cannot know what God is in himself by our natural power, he says, we proceed chiefly by learning how God is not: not bodily, not composed, not mutable, not contained in a genus (ST I, q.~3, prologue). Negation does not erase the positive knowledge gained from effects. It purifies that knowledge of the mode in which perfections belong to creatures.

God is good, but not as one good object among others. God is wise, but not by receiving truths, reasoning toward conclusions, or adding knowledge. God is present, but not by occupying coordinates. God is one, but not as one counted unit. Analogical language is neither literal in an identical creaturely sense nor merely decorative. It speaks truly because effects resemble their cause; it fails to contain because the mode of every divine perfection infinitely exceeds the mode from which the creaturely word was learned (ST I, q.~13, aa.~4--5).

\subsection{The luminous night}

The anonymous late-fifth- or early-sixth-century theologian known as Pseudo-Dionysius makes this purification into an ascent. The mind moves through sensible things and intelligible names, affirming every perfection of the universal cause; then it relinquishes the thought that even its highest affirmation captures God as he is. In the \emph{Mystical Theology}, Moses enters the darkness beyond seeing and being seen. The darkness signifies neither evil nor divine absence. It is the dazzling failure of finite sight before an excess of light (\emph{Mystical Theology} I.1--3; II.1).

This ``unknowing'' is not ignorance enthroned, a rejection of doctrine, or a technique for producing altered consciousness. Dionysian negation depends upon the affirmations it transcends. One must first know that God is cause, good, wise, living, and beyond creation before understanding why no creaturely form of those names circumscribes him. Silence is not the rival of truth. It is truth's reverence before the One it cannot turn into an object.

St.~Augustine narrates the existential form of the ascent. Turning inward, he passes beyond mutable bodily images and even beyond the mutable mind to an unchangeable Light above it---above because it made him, while he is below because he was made (\emph{Confessions} VII.10.16--11.17). He does not report calm mastery. The vision is flash-like; the mind knows, loves, trembles, and cannot remain by its own strength. In Book XI he holds together dread before the divine unlikeness and burning desire awakened by the likeness God has impressed upon the creature (XI.9.11).

The experience gives metaphysics a human temperature. To discover the uncreated source is not to stand above a successful argument. It is to discover that the power by which one argues, the truth one recognizes, and the being of the one who recognizes it are all received---while the One dimly known is not received at all.

\begin{studybox}{Apophatic realism}
The mind does not enter darkness because there is nothing to know. It enters because the reality known exceeds every finite manner of knowing. The right response is neither agnosticism nor conceptual possession, but a disciplined union of affirmation, negation, wonder, and desire.
\end{studybox}

At the summit of natural ascent, then, there is a strange conjunction: greater certainty that God is, greater inability to imagine what God is, greater intimacy of causal presence, and greater transcendence of essence. The mind has reached the shore of its natural power. It can know that the sea is there. It cannot, by staring harder at the water, discover the secret of the divine life.
